\conclusion

В работе был разработан и реализован подход к верификации соответствия TLA+
спецификации и исходного кода. В основе предложенного метода лежит
сопоставление трасс реального исполнения программы с допустимыми поведениями
формальной модели при помощи TLC.

В теоретической части были рассмотрены основы формальной спецификации в TLA+,
каноническая форма записи алгоритмов в виде машин состояний, а также формальная
постановка задачи валидации трасс как проверки непустоты пересечения множества
поведений спецификации и множества поведений, совместимых с наблюдаемой трассой
исполнения. Показано, что задача может быть сведена к ограниченной модельной
проверке посредством построения обёртки над исходной спецификацией.

В практической части была реализована библиотека инструментирования на языке C,
позволяющая логировать изменения абстрактных переменных спецификации и
связывать их с действиями TLA+. Разработана инфраструктура для автоматической
генерации трасс в формате NDJSON и их последующей проверки с использованием
TLC. Для спецификации алгоритма была создана ограниченная версия модели
(\texttt{TwoPhaseTrace}), обеспечивающая учёт ограничений, задаваемых трассой
исполнения.

Предложенная методика была апробирована на демонстрационном примере — алгоритме
двухфазной транзакции. Реализация менеджера транзакции и менеджеров ресурсов
была инструментирована и проверена на соответствие формальной спецификации. В
ходе валидации был обнаружен дефект реализации, связанный с некорректной
обработкой повторных сообщений \texttt{Prepared}, приводящий к преждевременной
фиксации транзакции. Обнаружение ошибки средствами трассовой верификации
подтверждает практическую ценность разработанного подхода.

Основным преимуществом предложенного метода является его применимость к
существующему коду без необходимости полной формальной верификации реализации
на уровне исходного или машинного кода. Подход позволяет выявлять логические
ошибки проектирования и расхождения между моделью и реализацией при
относительно невысокой стоимости внедрения. При этом степень детализации
трассировки может варьироваться, что обеспечивает баланс между точностью
проверки и вычислительными затратами.

К ограничениям метода относится его неполнота: анализируется лишь подмножество
поведений, соответствующее наблюдаемым трассам. Тем не менее, при достаточном
покрытии сценариев исполнения данный подход существенно повышает уверенность в
корректности реализации.

Таким образом, в работе показано, что интеграция формальной спецификации TLA+ с
трассовой валидацией исполнения представляет собой эффективный и практически
применимый способ повышения надёжности распределённых программных систем.
